\subsection{The terminal class period and the Bessel integral}
\label{sec:actual-theta-period-comparison}

There are two different period comparisons. The first identifies the
terminal pairing of the integral theta chain with an algebraic class
sum. The second unfolds the scalar Poincar\'e series into an integral
of Bessel periods. We prove both comparisons with their original
normalizations. An additional comparison between the two periods is
still required. The original source preserves linear relations in that
comparison; it does not identify their normalizing factors by declaring
them to be one.\cite{theta-period-source}

\paragraph{The finite terminal pairing.}
Keep an admitted datum of Theorem~\ref{thm:actual-theta-chain}, its
coefficient field $K$, embedding $\sigma:K\hookrightarrow\mathbb C$,
word prime $T$, original prime element $\pi_T$, and all its characters
and point representatives. Write $\mathcal C_i$ for its stage-$i$
classes, $w_i(c)=|\operatorname{Aut}(c)|$, and $V_i(K)$ for its
polynomial coefficient module. With the original kernels, set
\[
 x_0=\varphi,\qquad
 x_{i+1}(c')=\sum_{c\in\mathcal C_i}
                 \frac{K_i(c,c')x_i(c)}{w_i(c)}.
 \tag{TP1}
\]
The notation is independent of the indexing convention for the
single-step maps in Theorem~\ref{thm:actual-theta-chain}.

Let $\mathcal C'$ be the original terminal row classes. A representative
$d$ has rank $N_m-1$. Its prescribed lift $\widetilde d$ has rank
$N_m$, obtained by adjoining the original line with parameter $\pi_T$;
its frame is transported by the original change of representatives.
Write $x_m(\widetilde d)$ for the resulting polynomial value. Let
$B_{\rm pol}$ be the original normalized bilinear polynomial pairing,
and let $\mathcal B_T:V_m(K)\to V'(K)$ be the original branching map.
In particular this is a bilinear, not a Hermitian, pairing.

Here is an explicit description fixing the branching scalar. At an
embedding, let $\lambda$ and $\mu$ be the interlacing residual
partitions of ranks $N=n+1$ and $n$, and retain their common determinant
shift. Put $a_i=\mu_i-\lambda_{i+1}$ and
$b_i=\lambda_i-\mu_i$. The previously proved integral branching
formula, with its original normalization\cite{theta-period-branching}, uses the polynomial
$\xi_{\lambda,\mu}$ formed from the initial-column minors, with
exponents $a_i,b_i$, and the constant
\[
 c_{\lambda,\mu}=
 \prod_{i=1}^n\space
 \prod_{v\in [\lambda_{i+1},\,\min(\lambda_i,\lambda_{i+1}+
                                     \sum_j a_j)]\cap\mathbb Z
                       \setminus\{\mu_i\}}
 \prod_{j=1}^n\bigl((v-i)-(\mu_j-j)\bigr).
\]
Let $\Pi_\mu$ be the projection specified by the inverse Gram matrix
for $B_{\rm pol}$, and $[z^b]$ the coefficient of the indicated
last-row monomial, after discarding monomials involving the last
column. The branching map is the tensor product of
\[
 \tau(\overline{\pi_T})^{-n_{\lambda,\mu}}
 c_{\lambda,\mu}
 \frac{B_{\rm Fock}(\xi_{\lambda,\mu},\xi_{\lambda,\mu})}
      {B_{\rm Fock}(H_\lambda,H_\lambda)}
                 \Pi_\mu[z^b],
 \tag{TP2}
\]
where $H_\lambda$ is the standard extremal polynomial and
$n_{\lambda,\mu}$ is the difference of the sums of the original
weights, including their determinant shifts. This is precisely the
previously proved integral branching normalization. All powers of
$\overline{\pi_T}$ and all factorials in the Fock pairings remain.

For an algebraic row $z$, define
\[
 \mathcal P_K(\varphi,z)=
 \sum_{d\in\mathcal C'}\frac{
   B_{\rm pol}\bigl(\mathcal B_T x_m(\widetilde d),z(d)\bigr)}
                              {|\operatorname{Aut}(d)|}.
 \tag{TP3}
\]
The complex terminal period $\mathcal P_{\mathbb C}$ is the original
one: compose the original complex chain, evaluate at the same lifted
row points, apply its normalized lowering map, and pair with the
original row using the original class masses.

\begin{theorem}[Exact algebraic terminal period]
\label{thm:actual-algebraic-terminal-period}
For every original algebraic column $\varphi$ and row $z$,
\[
 \mathcal P_{\mathbb C}(\operatorname{ev}_\sigma\varphi,
                       \operatorname{ev}_\sigma z)
                         =\sigma\bigl(\mathcal P_K(\varphi,z)\bigr).
 \tag{TP4}
\]
The constant is one. The identity holds on the original twisted
initial source and on every retained source subspace. It introduces
no new finite denominator; the branching and pairing denominators
in (TP2)--(TP3) have not been removed.
\end{theorem}

\begin{proof}
The actual local-chain theorem identifies each point kernel, including
its root phases, with the corresponding polynomial kernel. Evaluation
commutes with a finite sum, with multiplication by a coefficient, and
with $w_i(c)^{-1}$. Induction in (TP1) therefore identifies the complex
chain at every stage with evaluation of $x_i$.

At the last two steps, the original coupled class space only changes
the representatives used for this sum. Its underlying decorated
lattice and automorphism group are the same as in the original
stage. In general, if a class is written by an intermediate choice
on which its automorphism group $\Gamma$ acts, the sum of the
weights $1/|\operatorname{Stab}_\Gamma(y)|$ over the choice orbits
equals $|Y|/|\Gamma|$. This is the orbit--stabilizer formula applied
to each orbit. Applying it to the original intermediate classes
recovers exactly the two coupled kernels, with no extra class count.
The final lift is the prescribed $\widetilde d$ because its
coefficient element and its conjugate satisfy
$\pi_T\overline{\pi_T}=T$ and avoid every retained auxiliary level.
Thus the same local lattices and residue functionals are used on
both sides.\cite{theta-period-source}

The normalized lowering map is (TP2) in the original frame. To check
its compatibility with evaluation, extract one polynomial coefficient
at a time, then apply the inverse Gram projection. Field embeddings
commute with both finite operations and with the displayed nonzero
Gram denominators. The original frame change precedes lowering on
both sides. Its equivariance follows from interlacing and the common
determinant shift, already proved for the exact branching map.

The two terminal polynomial values transform by opposite character
factors. The first carries the inverse factor and the second the
factor itself; $B_{\rm pol}$ is invariant under the corresponding
dual matrix actions. Their product is therefore invariant, with
its phase exactly cancelled. The point/phase equivalence preserves
the stabilizer orders. The original terminal period is consequently
the image under $\sigma$ of precisely (TP3). This proves (TP4).
Precomposing the identity with the already specified initial twist
or restricting its argument proves the final assertions.
\end{proof}

\paragraph{The analytic period, with its Haar factor.}
Fix one original global word $u$ in the admitted arithmetic gap regime,
where the original scalar Poincar\'e series belongs to the specified
$L^2$ space. Use its rational flag and its specified
nonisotropic vector to write the Bessel stabilizer as
$S_u=N_u\rtimes H_u$. Its real and finite parts are
$S_{u,\infty}$ and $S_{u,f}$; $H_u$ is the original anisotropic
unitary group. These are subgroups of the original ambient group,
not abstract groups substituted for them. Keep the rational subgroup
$S_u(F)$, compact open $K_2$, finite representative $x_u$, the
original scalar seed $\Phi_{u,T}$, and its periodization $P_uT$.

Let $\mu_f$ be the original finite Haar measure and $v_K=\mu_f(K_2)$.
Normalize an auxiliary Haar measure $\nu_{u,f}$ on $S_{u,f}$ by
\[
 \nu_{u,f}(S_{u,f}\cap x_uK_2x_u^{-1})=1.
\]
The real stabilizer measure is the prescribed product measure
$\mu_{u,S}=\mu_{u,N}\mu_{u,H}$. The original Bessel Haar measure
$\mu_{u,B}$ has a unique conversion factor $\beta_u>0$:
\[
 \mu_{u,B}=\beta_u(\mu_{u,S}\times\nu_{u,f})
 \tag{TP5}
\]
under the original rational-model identification. This is a specified
ratio of Haar measures, not a free normalization. It can equivalently
be computed by evaluating the two measures on any fixed relatively
compact set of positive measure.

Let $\nu_{u,Q}$ be the original real quotient measure and $W_u(y)$
the original vector-valued seed. For a bounded continuous automorphic
row $f$ with right character $\overline{\chi_2}$, put
\[
 \mathcal I_u(T,f;y)=
 \int_{S_u(F)\backslash S_u(\mathbb A)}
  \Psi_u(s)\,T(W_u(y))(h_s)\,
                    f\bigl(s(y,x_u)\bigr)\,d\mu_{u,B}(s).
 \tag{TP6}
\]
Here $s=n_sh_s$ in the fixed semidirect product, $T$ is read using
the original column dictionary, and $\Psi_u$ is the inverse Bessel
character: in the original coordinates it is the product of the
prescribed infinite and finite unipotent characters. The integral
uses the original quotient measure, retaining the finite class
masses. This is the original Bessel period, with its radial variable
$y$ still present.

\begin{theorem}[Normalized Poincar\'e--Bessel pairing]
\label{thm:actual-poincare-bessel-pairing}
Let $\phi$ be the $L^2$ class of the row $f$. With the Hilbert
inner product conjugate-linear in its first argument, one has
\[
 \left\langle\overline\phi,\,c_{0,u}^{-1/2}P_uT\right\rangle
   =\frac{v_K}{\beta_u\sqrt{c_{0,u}}}
             \int\mathcal I_u(T,f;y)\,d\nu_{u,Q}(y),
 \tag{TP7}
\]
where the original positive normalizer is
\[
 c_{0,u}=\frac{v_K\,\mu_{u,N}(\mathcal F_{u,N})
                         \mu_{u,H}(H_{u,\infty})}{\dim V_u}.
\]
Here $\mathcal F_{u,N}$ is the original real fundamental domain for
the rational unipotent subgroup whose finite component preserves
$x_uK_2x_u^{-1}$. In particular its volume has not been replaced by
one.
All integrals in this identity are absolutely convergent. The factor
in (TP7) is fixed before choosing the test column or row and is
strictly positive. No uniform estimate for $\beta_u$ is asserted.
\end{theorem}

\begin{proof}
First consider a nonnegative function supported on $S_{u,f}x_uK_2$.
For $g$ in this set, the set of $s\in S_{u,f}$ with
$s^{-1}g\in x_uK_2$ is a left coset of
$S_{u,f}\cap x_uK_2x_u^{-1}$, so its $\nu_{u,f}$-mass is one.
Fubini and left invariance of $\mu_f$ give
\[
 \int F(g)\,d\mu_f(g)=
 \int_{S_{u,f}}\int_{K_2}F(sx_uk)\,d\mu_f(k)\,d\nu_{u,f}(s).
 \tag{TP8}
\]
Combine this with the original real Weil formula and (TP5).
The resulting adelic disintegration factor is $\beta_u^{-1}$.

Choose a measurable fundamental domain for the discrete rational
subgroup $S_u(F)$. Folding the preceding disintegration to that
domain is justified first for nonnegative functions by Tonelli.
The rational translates of its indicator add to one almost
everywhere. For a right-$K_2$-invariant function, the remaining
integration over $K_2$ gives exactly $v_K$. Thus the folded factor
is $v_K/\beta_u$. No stabilizer or level volume has been discarded.

For the seed, the original character law gives, on its support,
\[
 \Phi_{u,T}\bigl(s(y,x_u)\bigr)
                       =\Psi_u(s)T(W_u(y))(h_s).
\]
This is the original column evaluation, read at the rational flag
and the same finite representative. The finite-level covariance
of $f$ is the inverse of that of the seed; their product is
right-$K_2$ invariant. The original compact anisotropic quotient
and the finite-dimensional column dictionary give a bound
$|\Phi_{u,T}(s(y,x_u))|\le C_T\|W_u(y)\|$ on the compact
rational Bessel quotient. The proved seed decay gives
$\int\|W_u(y)\|\,d\nu_{u,Q}(y)<\infty$.
The row $f$ is bounded. Hence the same folded formula proves
absolute convergence for the seed and for $f\Phi_{u,T}$.

The original absolute Poincar\'e convergence then permits unfolding
the periodized integral:
\[
 \langle\overline\phi,P_uT\rangle
  =\int_{S_u(F)\backslash G(\mathbb A)}f(g)\Phi_{u,T}(g)\,dg
  =\frac{v_K}{\beta_u}
                        \int\mathcal I_u(T,f;y)\,d\nu_{u,Q}(y).
\]
The first equality has $f$, rather than $\overline f$, because the
first argument is $\overline\phi$. Division by the original
positive square root proves (TP7).
\end{proof}

The finite comparison (TP4) and the analytic comparison (TP7) remove
two actual period-identification obligations. They do not identify
the radial Bessel integral with (TP3). The remaining comparison must
retain the actual smooth row realization, its restrictions and the
see-saw identities for the intermediate theta lifts. In the source,
two-term linear relations are transported first; a nonzero Poincar\'e
pairing then supplies nonzero proportionality factors. These factors
may depend on the row tuple. The original period normalizer $c_{0,u}$
does not by itself determine them. Thus neither W06 nor the final
source and target trace bounds are concluded here.
